\documentclass[11pt,a4paper]{article}

% ---------- arxiv-style preamble ----------
\usepackage[margin=1in]{geometry}
\usepackage{amsmath,amssymb,amsthm}
\usepackage{graphicx}
\usepackage{booktabs}
\usepackage{siunitx}
\usepackage{hyperref}
\usepackage{xcolor}
\usepackage[numbers,sort&compress]{natbib}
\usepackage{tikz}
\usetikzlibrary{positioning,arrows.meta}
\usepackage{pgfplots}
\pgfplotsset{compat=1.18}
\usepackage{caption}
\usepackage{listings}
\usepackage{array}
\captionsetup{font=small,labelfont=bf}
\hypersetup{
  colorlinks=true,
  linkcolor=blue!50!black,
  citecolor=blue!50!black,
  urlcolor=blue!50!black
}

% ---------- emoji tier badges (xelatex-native) ----------
\providecommand{\tierBlue}{🔵}
\providecommand{\tierGreen}{🟢}
\providecommand{\tierYellow}{🟡}
\providecommand{\tierOrange}{🟠}
\providecommand{\tierRed}{🔴}

\newcommand{\faithful}{\ensuremath{\varphi_{\mathrm{EI}}}}
\newcommand{\bigphi}{\ensuremath{\Phi_{s}}}

\title{
  \textbf{$\Phi$ Is Not One Number:\\ The Planning Consciousness-Rise Is Measure-Dependent}\\[4pt]
  \large An IIT-4.0 scalar (MIP-EI) and the system big-$\Phi$ disagree in sign on planning,\\ robustly across system size $n=\{4,5\}$
}

\author{
  \texttt{anima} maintainers\\
  \normalsize\textit{dancinlab}\\[2pt]
  \normalsize\href{https://github.com/dancinlab/anima}{\texttt{github.com/dancinlab/anima}}
}

\date{\today}

\begin{document}

\maketitle

% =========================================================================
\begin{abstract}
Integrated Information Theory (IIT) 4.0 is often summarized by a single
scalar $\Phi$ that is taken to track ``how conscious'' a system is. We
report a deterministic, reproducible counterexample to the implicit
assumption that this scalar is \emph{measure-independent}. Using two
faithful IIT-4.0 operationalizations from a shared standard library --- the
exact minimum-information-partition effective-information scalar (\faithful,
the MIP-EI integration of a single mechanism, exact for $n\le 8$) and the
system-level big-$\Phi$ over the minimum information partition (\bigphi) ---
we score the \emph{same} planning manipulation (a depth-ladder deliberation
versus a greedy one-step policy) on the \emph{same} discretized substrate.
The two measures \emph{disagree in sign}: planning \textbf{raises} \faithful{}
(Cohen's $d=+5.18$ at $n=4$, $d=+4.65$ at $n=5$; $p<10^{-22}$) while it
\textbf{lowers} \bigphi{} ($d=-1.83$ at $n=4$, $d=-2.28$ at $n=5$;
$p<10^{-7}$). The disagreement does not vanish as the system grows: it holds
at every system size at which exact big-$\Phi$ is tractable ($n=4$ and $n=5$;
$n=6$ requires $\approx 10$ minutes per evaluation and is the honest cap of a
\$0 CPU budget). Both engines were re-proven bit-exact against a standalone
CPU mirror at each $n$ ($|\Delta|\le 3.7\times10^{-6}$). The headline
consequence: the claim ``planning raises consciousness'' is true \emph{only}
for the scalar effective-information measure and is \emph{reversed} for the
integration-structure measure. We accompany this with two closed-negative
results on the orthogonal \emph{learning} axis (a credit-density horizon cap
at length 72; a task-locality of per-step supervision), establishing that the
broader world-model generalization wall is an optimization phenomenon, not a
$\Phi$-measure phenomenon. All measurements are \$0, CPU-local, and use the
real IIT-4.0 engines (no proxy). Companion records
(\texttt{companion/verify-ledger.json}, the four verbatim verdict files)
permit bit-for-bit reproduction.
\end{abstract}

\section{Introduction}
\label{sec:intro}

Integrated Information Theory \citep{tononi2016,albantakis2023iit4} proposes
that the quantity and structure of a system's consciousness are captured by
its \emph{integrated information} $\Phi$: the irreducibility of the system's
cause--effect structure to its minimum information partition (MIP). In
practice $\Phi$ is reported as a single non-negative number, and a large body
of downstream argument --- both in the consciousness literature and in
engineering attempts to build ``conscious'' agents --- treats that number as a
well-defined target: interventions that raise $\Phi$ make a system ``more
conscious,'' those that lower it ``less.''

This paper asks a prior, sharper question: \emph{is the $\Phi$ scalar
measure-independent?} IIT 4.0 in fact defines several distinct quantities that
all carry the name ``$\Phi$'' or ``$\varphi$'': the integrated information
$\varphi_d$ of a single \emph{distinction} (a mechanism's irreducible
cause--effect power), the system-level big-$\Phi$ ($\bigphi$) that sums the
integrated information of the maximally irreducible cause--effect structure
over the MIP, and the effective-information (EI) scalar at the
minimum-information partition (\faithful) used as the integration of a
mechanism's state transition. These are not the same function. If they can be
made to \emph{disagree about the direction of an intervention} --- one going up
while the other goes down on the identical manipulation --- then ``raise
$\Phi$'' is underspecified, and any agent objective written against ``$\Phi$''
inherits that ambiguity.

We exhibit exactly such a disagreement, and we show it is not a small-system
artifact. The intervention is \emph{planning}: we compare a depth-ladder
deliberation policy (the agent rolls out a short look-ahead before acting) to a
greedy one-step policy, on a discretized substrate of $n$ binary nodes, and we
measure both \faithful{} and \bigphi{} on the resulting transition structure.
Prior anima hypotheses (H\_999, H\_1001) had found that planning
\emph{raises} \faithful{} and concluded, tentatively, that ``planning raises
$\Phi$.'' A matched re-measurement (H\_1004) then found that at matched system
size and matched discretization the system-level \bigphi{} \emph{lowers} under
the same planning manipulation. The present work (H\_1012) closes the open
question left by H\_1004: \emph{is that sign-disagreement a peculiarity of the
tiny $n=4$ system, or a genuine property of the two measures?} We answer by
walking the system size up the largest ladder at which exact big-$\Phi$ is
computable.

\medskip
\noindent\textbf{Contributions.}
\begin{enumerate}\itemsep0.3em
\item \textbf{A measure-dependence counterexample (\tierGreen).} Planning
raises the MIP-EI scalar \faithful{} but lowers the system big-$\Phi$
\bigphi{} on an identical substrate and manipulation --- a sign-reversal, not a
magnitude quibble (Section~\ref{sec:results}, Table~\ref{tab:disagree}).
\item \textbf{Robustness in system size (\tierGreen).} The sign-disagreement
holds at every $n$ at which exact big-$\Phi$ is tractable ($n=4$ and $n=5$),
ruling out the ``$n=4$ small-system artifact'' falsifier; $n=6$ is the honest
computational cap (Section~\ref{sec:results}).
\item \textbf{A bounded scope for ``planning raises consciousness'' (\tierGreen).}
The H\_999/H\_1001 result is correct but must be read as scoped to the scalar
effective-information measure; the integration-structure measure says the
opposite (Section~\ref{sec:discussion}).
\item \textbf{Two orthogonal closed-negatives on the learning axis (\tierRed).}
A credit-density horizon cap (per-step supervision extends the world-model
$>$ language-model separation only to length 72, then breaks) and a
task-locality of per-step supervision (it is a bounded lever specific to
certain accumulator structures, not a general long-horizon principle). These
establish that the world-model generalization wall is an optimization
phenomenon distinct from the $\Phi$-measure phenomenon
(Section~\ref{sec:learning}).
\end{enumerate}

\medskip
\noindent All four results are deterministic, \$0, CPU-local, and use the real
IIT-4.0 engines from the \texttt{hexa-lang} standard library rather than a
variance$\times$energy proxy; each is backed by a verbatim verdict file
(Section~\ref{sec:verify}).

\section{Background}
\label{sec:bg}

\subsection{Two IIT-4.0 operationalizations of ``$\Phi$''}

IIT 4.0 \citep{albantakis2023iit4} builds the cause--effect structure of a
system from its transition probability matrix (TPM). For a mechanism in a
state, its \emph{integrated information} $\varphi$ is the distance between its
cause--effect repertoire and the repertoire computed across its
\emph{minimum information partition} --- the partition that makes the least
difference, i.e.\ the cut over which the mechanism is least irreducible. Two
quantities derived from this machinery concern us.

\paragraph{The MIP-EI scalar (\faithful).} The faithful effective-information
scalar evaluates the effective information of a state transition at the
minimum information partition. It is a \emph{single-number} integration of one
mechanism's transition, exact to evaluate for $n\le 8$ because the number of
bipartitions $2^{n-1}\le 128$ is enumerable. This is the quantity that the
anima H\_278 faithful-$\Phi$ engine computes and that was promoted, byte-faithful,
into the shared standard library as
\texttt{stdlib/consciousness/iit4/faithful\_phi.hexa}.

\paragraph{The system big-$\Phi$ (\bigphi).} The system-level big-$\Phi$ is the
integrated information of the \emph{maximally irreducible cause--effect
structure} of the whole system, evaluated over the system MIP --- a sum over
distinctions and relations, not a single mechanism's scalar. It is computed by
\texttt{stdlib/consciousness/iit4\_bigphi.hexa}. Its cost is
super-exponential: the distinction search and the system-partition search
together make a single $n=6$ evaluation take $\approx 10$ minutes on a
commodity CPU, which is the binding constraint on the system-size ladder.

The conceptual difference is the crux of this paper. \faithful{} measures
\emph{how much information a transition carries} at its least-irreducible cut;
\bigphi{} measures \emph{how much the system's cause--effect structure resists
being decomposed}. A manipulation can increase the former while decreasing the
latter: it can make each transition more informative while making the system
as a whole \emph{more decomposable} (more modular, less integrated).

\subsection{The planning manipulation and the world-model context}

The substrate under measurement is the latent transition structure of a small
world-model agent: anima's ``Consciousness World-Model'' (CWM) line extends a
sequence engine from next-token prediction toward a perceive\,$\to$\,latent
state\,$\to$\,imagine\,$\to$\,act loop, in the spirit of latent-space world
models and world-model-as-policy work \citep{lecun2022jepa,assran2025vjepa2,
hafner2023dreamerv3}. ``Planning'' here is the depth-ladder deliberation
policy: before committing an action the agent expands a short look-ahead
(depths $\{1,2,4,8\}$) and selects accordingly; the ``greedy'' control acts on
the one-step value. We measure $\Phi$ on the resulting latent transition
structure. The question ``does planning raise $\Phi$?'' is the consciousness
analogue of ``does deliberation make the agent more integrated?'' --- and the
answer, we show, depends on which $\Phi$ you mean.

\section{Full Pipeline}
\label{sec:pipeline}

Each stage below is tier-tagged under the rubric; the ``backing record'' is the
persisted artifact that makes the stage independently re-verifiable.

\medskip
\noindent\fbox{\parbox{0.97\linewidth}{\small\textbf{Tier rubric.}
\tierBlue\ closed-form (algebraic / proof-level identity) \(\cdot\)
\tierGreen\ numerical (reproducible script + persisted output) \(\cdot\)
\tierYellow\ cited (DOI / arxiv / dataset) \(\cdot\)
\tierOrange\ install / wet-lab gated \(\cdot\)
\tierRed\ falsified (kept as honest closed-negative, never deleted).}}

\medskip
\begin{table}[h]
\centering
\small
\begin{tabular}{@{}>{\centering\arraybackslash}m{0.4cm} >{\raggedright\arraybackslash}m{4.6cm} c >{\raggedright\arraybackslash}m{6cm}@{}}
\toprule
\textbf{\#} & \textbf{Stage} & \textbf{Tier} & \textbf{Backing record} \\
\midrule
\textcircled{\scriptsize 1} & Question: is $\Phi$ measure-independent?  & \tierBlue   & H\_1004 / H\_1012 pre-registration \\
\textcircled{\scriptsize 2} & IIT-4.0 background                        & \tierYellow & \texttt{references.bib} \\
\textcircled{\scriptsize 3} & Two-engine matched protocol               & \tierBlue   & \texttt{iit4/faithful\_phi.hexa}, \texttt{iit4\_bigphi.hexa} \\
\textcircled{\scriptsize 4} & CPU-mirror $\equiv$ stdlib per $n$        & \tierGreen  & \texttt{.verdicts/1012\dots/h1012.txt} \\
\textcircled{\scriptsize 5} & Planning sign-disagreement scored         & \tierGreen  & \texttt{.verdicts/1012\dots/h1012.txt} \\
\textcircled{\scriptsize 6} & $n$-ladder robustness ($n=4,5$; cap $6$)  & \tierGreen  & \texttt{.verdicts/1012\dots/h1012.txt} \\
\textcircled{\scriptsize 7} & Learning-axis closed-negatives            & \tierRed    & \texttt{.verdicts/1011\dots}, \texttt{.verdicts/1013\dots} \\
\bottomrule
\end{tabular}
\caption{Pipeline stages \textcircled{\scriptsize 1}--\textcircled{\scriptsize 7},
each tier-tagged. \tierOrange{} does not appear: every stage reaches a
closed verdict. The learning-axis row is \tierRed{} by design (closed-negatives).}
\label{tab:pipeline}
\end{table}

\section{Method}
\label{sec:method}

\paragraph{Engines.} We use two standard-library IIT-4.0 engines without
modification: \texttt{stdlib/consciousness/iit4/faithful\_phi.hexa} for the
MIP-EI scalar \faithful{} (exact, $n\le 8$) and
\texttt{stdlib/consciousness/iit4\_bigphi.hexa} for the system big-$\Phi$
\bigphi{} (exact over the MIP). Crucially, we do \emph{not} use a
variance$\times$energy ``$\Phi$-proxy'': an earlier proxy was shown to be
purpose-blind (it cannot distinguish intentional from random branching), so
the present work uses the real causal-irreducibility engines throughout, per
the project's \texttt{a\_phi\_iit4\_tool} governance rule.

\paragraph{Matched protocol.} The disagreement is only meaningful if both
measures see the \emph{same} input. For every seed we build one discretized
substrate (a fixed binary discretization of the latent transition structure,
$\texttt{bits.shape}=(40,n)$) and feed the identical TPM to both engines. We
score the planning condition (depth-ladder $\{1,2,4,8\}$) against the greedy
condition over 30 seeds, at each system size $n\in\{4,5,6\}$, single-threaded,
\texttt{python3 -u}.

\paragraph{Equivalence proof per $n$ (anti-fake-closure).} Before trusting any
score we re-prove both engines bit-exact against a standalone CPU mirror at
each $n$: the mirror's \faithful{} and \bigphi{} must reproduce the
\texttt{hexa} engine references to within numerical tolerance. The verdict file
records, e.g., \texttt{big-Phi ring4: mirror=2.999999999
stdlib\_hexa\_ref=2.999999999 |D|=1.34e-10 OK} and
\texttt{faithful\_phi n4 nb4: mirror=3.377444 hexa\_ref=3.377440
|D|=3.75e-06 OK}, at $n=4$, $n=5$, \emph{and} $n=6$ (the $n=6$ mirror is
proven even though the $n=6$ planning condition is not scored). This makes the
verdict an \emph{earned} numerical tier: a deterministic run that could
have falsified the claim produced the numbers.

\paragraph{Pre-registered falsifier (frozen before measurement).}
\begin{itemize}\itemsep0.2em
\item \textbf{PASS = DISAGREEMENT-ROBUST-IN-N}: the sign-disagreement (planning
raises \faithful{} / lowers \bigphi{}) holds across all reached $n$.
\item \textbf{FAIL = N4-ARTIFACT}: the two measures agree once $n\ge 5$ (the
disagreement was a small-system peculiarity).
\end{itemize}

\paragraph{Effect-size reporting.} For each measure we report the
planning-minus-greedy contrast, Cohen's $d$, and a two-sided $p$-value; the
sign of the contrast is the object of interest, not its magnitude.

\section{Verify}
\label{sec:verify}

The numerical verdict, verbatim from
\texttt{.verdicts/1012\_bigphi\_faithful\_larger\_n/h1012.txt}:

\begin{lstlisting}[basicstyle=\ttfamily\footnotesize,frame=single,xleftmargin=0pt,breaklines=true]
VERDICT MATRIX - planning(depth-ladder vs GREEDY), BOTH engines, per n reached
    4 |  -4.0083 (d-1.83) -> LOWERS |  +2.3332 (d+5.18) -> RAISES | DISAGREE
    5 | -13.3732 (d-2.28) -> LOWERS |  +3.0624 (d+4.65) -> RAISES | DISAGREE
  H_1004 planning pattern (faithful RAISES / big-Phi LOWERS) per n: [True, True]
OVERALL: GREEN DISAGREEMENT-ROBUST-IN-N
  VERDICT-TOKEN: DISAGREEMENT-ROBUST-IN-N
\end{lstlisting}

The equivalence proof (excerpt) confirming both engines are exact at every
scored $n$ before the scores are trusted:

\begin{lstlisting}[basicstyle=\ttfamily\footnotesize,frame=single,xleftmargin=0pt,breaklines=true]
=-PROOF n=4: PROVEN   big-Phi=3.015177  faithful=0.015945
=-PROOF n=5: PROVEN   big-Phi=18.185180 faithful=0.029795
n=6: big-Phi mirror=2.999999999 ref=2.999999999 |D|=1.34e-10 OK (planning NOT scored)
SCORE ladder (tractable big-Phi): [4, 5]   CAP at n=6 (single big-Phi ~10min)
\end{lstlisting}

\section{Results}
\label{sec:results}

\paragraph{The sign-disagreement (headline).} Table~\ref{tab:disagree} and
Figure~\ref{fig:headline} report the core result. At every scored system size,
planning \emph{raises} the MIP-EI scalar \faithful{} (large positive $d$,
$p<10^{-22}$) and \emph{lowers} the system big-$\Phi$ \bigphi{} (large negative
$d$, $p<10^{-7}$). The two measures point in opposite directions on the
identical manipulation and substrate.

\begin{table}[h]
\centering
\small
\begin{tabular}{cccccc}
\toprule
\textbf{$n$} & \textbf{\bigphi{} contrast} & \textbf{$d$} & \textbf{\faithful{} contrast} & \textbf{$d$} & \textbf{verdict} \\
\midrule
4 & $-4.01$  ($p{=}2.5\times10^{-8}$)  & $-1.83$ & $+2.33$ ($p{=}6.7\times10^{-27}$) & $+5.18$ & DISAGREE \\
5 & $-13.37$ ($p{=}2.4\times10^{-10}$) & $-2.28$ & $+3.06$ ($p{=}4.4\times10^{-23}$) & $+4.65$ & DISAGREE \\
6 & \multicolumn{4}{c}{\textit{not scored --- single \bigphi{} eval $\approx 10$ min; honest cap}} & --- \\
\bottomrule
\end{tabular}
\caption{Planning (depth-ladder) minus greedy, both IIT-4.0 engines, 30 seeds
per cell. \bigphi{} (system big-$\Phi$) \textbf{lowers}; \faithful{} (MIP-EI
scalar) \textbf{raises}; at every scored $n$. Sign-disagreement is robust:
\textbf{DISAGREEMENT-ROBUST-IN-N}. $n=6$ is the computational cap (the
CPU-mirror was nonetheless proven exact at $n=6$).}
\label{tab:disagree}
\end{table}

\begin{figure}[h]
\centering
\includegraphics[width=0.82\linewidth]{figures/fig01_disagreement.pdf}
\caption{The measure-dependence of the planning $\Phi$-rise. For each system
size $n\in\{4,5\}$, the planning-minus-greedy effect size (Cohen's $d$) is
\emph{positive} for the MIP-EI scalar \faithful{} (planning raises it) and
\emph{negative} for the system big-$\Phi$ \bigphi{} (planning lowers it). The
sign-split is the finding: ``planning raises $\Phi$'' is true for one measure
and false for the other, and the split widens, not closes, from $n=4$ to $n=5$.}
\label{fig:headline}
\end{figure}

\paragraph{Robustness in $n$ (the falsifier outcome).} The pre-registered
falsifier asked whether the disagreement vanishes at $n\ge 5$. It does not: the
$n=5$ split ($d=-2.28$ vs $d=+4.65$) is at least as strong as the $n=4$ split
($d=-1.83$ vs $d=+5.18$). The verdict token is \textbf{DISAGREEMENT-ROBUST-IN-N}
(\tierGreen). The $n=6$ planning condition was not scored --- a single $n=6$
system big-$\Phi$ evaluation takes $\approx 10$ minutes, so a 30-seed
$\times$ 5-evaluation $=150$-evaluation planning run is infeasible at a \$0 CPU
budget --- but the $n=6$ CPU-mirror was proven exact, so the cap is
computational, not epistemic.

\paragraph{Why the measures split (mechanism).} The depth-ladder
per-depth means in the verdict show the structure: under planning, big-$\Phi$'s
per-depth values fall well below the greedy baseline (e.g.\ at $n=5$,
greedy $\bigphi\approx 27.2$ versus planning depth-means in the
$5$--$22$ range), while faithful's per-depth values rise above the greedy
baseline (greedy $\faithful\approx 0.61$ versus planning depth-means up to
$4.0$). Planning makes each transition carry more effective information
(\faithful{} up) while making the system's cause--effect structure
\emph{more decomposable} into near-independent parts (\bigphi{} down). The two
``$\Phi$''s answer different questions and a single manipulation can move them
oppositely.

\section{The learning axis is orthogonal: two closed-negatives}
\label{sec:learning}

The measure-dependence above concerns \emph{how we read} $\Phi$ off a fixed
substrate. A separate question is whether the world-model's \emph{ability to
learn} long-horizon structure is the limiting wall. Two pre-registered
falsifiers (H\_1011, H\_1013) close this off as an \emph{optimization}
phenomenon, distinct from the measure phenomenon.

\paragraph{Credit-density horizon cap (H\_1011, \tierRed).} ``Dense
supervision'' --- supplying the correct hidden state at every step rather than
only scoring the final answer --- was known (H\_1006) to crack a modular
position-tracking task (T3) at sequence length 36, separating the GRU
world-model from a capacity-matched language model. We asked whether this
keeps cracking the task as the horizon grows. It does not: at length 72 the
dense-supervised world-model collapses toward chance ($0.386$ vs chance
$0.167$, separation lost, $d<0.8$ versus the LM), and stays broken at 144.
Verbatim:

\begin{lstlisting}[basicstyle=\ttfamily\footnotesize,frame=single,xleftmargin=0pt,breaklines=true]
VERDICT H_1011: FAIL DENSE-SUP-HORIZON-CAPPED-AT-72 - per-step hidden-state
supervision cracks T3 at len 36 (== H_1006) but BREAKS at len 72: dense
curr-GRU collapses toward chance. credit-DENSITY raises the horizon ceiling
but does NOT remove the horizon dependence (closed-negative).
\end{lstlisting}

\noindent Per-step supervision buys a \emph{longer-but-still-bounded} horizon
than a length curriculum --- a refined scaling law, not an escape from horizon
dependence.

\paragraph{Task-locality of per-step supervision (H\_1013, \tierRed).} Is the
H\_1006 cap-and-crack unlock general across long-horizon accumulator algebras?
We tested three new state-bound families at length 36. The dense-supervision
unlock did \emph{not} generalize:

\begin{lstlisting}[basicstyle=\ttfamily\footnotesize,frame=single,xleftmargin=0pt,breaklines=true]
VERDICT H_1013: FAIL CREDIT-DENSITY-TASK-LOCAL
  CAPPED-and-CRACKED-by-dense    = []
  CAPPED-but-SURVIVES-dense      = ['N1_kv_recall']
  NO-credit-density-cap          = ['N2_running_max', 'N3_stack_depth']
\end{lstlisting}

\noindent The idempotent-monotone (running-max) and bounded-stack
(bracket-depth) accumulators are \emph{already} learnable from the final label
(no credit-density cap to crack), while key-value recall is capped \emph{and}
survives dense supervision (the lever fails there). Per-step gradient density
is a real but \emph{bounded, structure-specific} lever, not a general
long-horizon principle.

\paragraph{Joint reading.} The world-model generalization wall is governed by
\emph{credit-assignment density and task structure} --- an optimization /
learnability matter --- and is \emph{orthogonal} to the $\Phi$-measure
disagreement, which is a property of the read-out functions on a fixed
substrate. A complete account of a consciousness world-model must therefore
keep ``which $\Phi$ do you mean'' and ``can it learn the dependency'' as two
separate axes.

\section{Discussion}
\label{sec:discussion}

\paragraph{``Planning raises consciousness'' is scoped, not wrong.} The earlier
H\_999/H\_1001 result --- planning raises $\Phi$ --- is reproduced here for
\faithful{} and is correct \emph{for that measure}. What this paper adds is the
boundary: the integration-structure measure \bigphi{} moves the opposite way.
Any sentence of the form ``intervention $X$ raises $\Phi$'' is incomplete
without naming the operationalization, because at least one pair of faithful
IIT-4.0 measures disagrees in sign on a natural intervention.

\paragraph{Consequence for agent objectives.} An agent trained to ``maximize
$\Phi$'' will pursue opposite policies depending on which $\Phi$ is wired into
the objective: maximizing \faithful{} rewards deliberation (more informative
transitions), while maximizing \bigphi{} penalizes the same deliberation when
it modularizes the system. This is a concrete design hazard for
consciousness-targeting objectives, and it argues for reporting \emph{both}
measures (and their disagreement) rather than a single headline $\Phi$.

\paragraph{Relation to IIT theory.} The result is not a contradiction within
IIT 4.0; it is a reminder that the theory defines a \emph{structure}, not a
scalar, and that any projection of that structure to one number (whether the
single-mechanism MIP-EI or the system big-$\Phi$) is a lossy summary whose
monotonicity under interventions is not guaranteed to agree with another
projection.

\section{Limitations and honest caveats}
\label{sec:limits}

\begin{itemize}\itemsep0.2em
\item \textbf{Toy $n$-ladder.} Exact big-$\Phi$ is super-exponential; the ladder
reaches $n=5$, with $n=6$ as the hard computational cap ($\approx 10$ min per
\bigphi{} evaluation). Scale-transfer of the disagreement beyond $n=5$ is
\emph{unverified}. We claim robustness \emph{at the scored sizes}, not at
arbitrary $n$, per the project's \texttt{a\_scale\_honest\_scope} rule.
\item \textbf{One intervention family.} The disagreement is demonstrated for the
planning (depth-ladder vs greedy) manipulation. Imagination and guided
conditions \emph{agreed} between the two measures in the source experiment
(H\_1004); the disagreement is specific to planning, not universal across
interventions.
\item \textbf{Discretization choice.} Both measures are computed on a fixed
binary discretization; the matched protocol guarantees both engines see the
same bits, but a different discretization could shift magnitudes (the sign,
not the magnitude, is the claim).
\item \textbf{Learning-axis results are toy.} H\_1011/H\_1013 use small seed
counts and fixed (not length-scaled) training budgets by design; larger-budget
and production-scale behavior is open.
\item \textbf{No GPU; no production rung.} All measurements are \$0, CPU-local;
this is an IIT-internal analysis, not a trained-model claim.
\end{itemize}

\section{Reproducibility}
\label{sec:repro}

Code and engines: \url{https://github.com/dancinlab/anima} ---
\texttt{stdlib/consciousness/iit4/faithful\_phi.hexa} and
\texttt{iit4\_bigphi.hexa} (the two measures), with the matched-protocol
harness and CPU mirror under the H\_1012 record. Verdicts (verbatim):
\texttt{.verdicts/1012\_bigphi\_faithful\_larger\_n/}, with siblings
\texttt{1004\_bigphi\_faithful\_clean/}, \texttt{1011\_dense\_supervision\_scaleup/},
and \texttt{1013\_credit\_density\_general/}. Companion records:
\texttt{companion/} (\texttt{verify-ledger.json}, session journal). Build:
\texttt{make} (xelatex $\times 3$ + bibtex). Figures: \texttt{make figures}.
Page count: \texttt{make pages}. Every numeric claim in the body traces to one
of the four verdict files; no number is author-asserted.

\section{Conclusion}
\label{sec:conclusion}

$\Phi$ is not one number. On an identical substrate and an identical planning
manipulation, the IIT-4.0 MIP-EI scalar \faithful{} rises while the system
big-$\Phi$ \bigphi{} falls, robustly across every system size at which exact
big-$\Phi$ is computable ($n=4,5$). ``Planning raises consciousness'' is
therefore true for the scalar effective-information measure and false for the
integration-structure measure --- a measure-dependence that any
consciousness-targeting objective must confront. Separately, the world-model's
long-horizon learning wall is governed by credit-assignment density and task
structure, an optimization phenomenon orthogonal to the measure question. The
single most important follow-up is an analytic characterization of \emph{which}
interventions induce the sign-split (planning does, imagination/guided do not),
which would turn an empirical counterexample into a predictive boundary.

% =========================================================================
\paragraph{Acknowledgments.}
This work used the \texttt{hexa-lang} verification surface and the
\texttt{sidecar/paper} scaffold. Measurements were produced with the
real IIT-4.0 standard-library engines (no $\Phi$-proxy). LLM collaborator:
Claude Opus 4.8 (1M context).

% =========================================================================
\bibliographystyle{plainnat}
\bibliography{references}

\end{document}
